% Standalone problem document, generated from the adjacent README.md.
% Compile with XeLaTeX. Mathematical content is maintained in Markdown.
\documentclass[11pt,a4paper]{article}
\usepackage[fontsize=10.5pt]{scrextend}
\usepackage[margin=22mm,headheight=15pt,headsep=9mm,footskip=11mm]{geometry}
\usepackage{fontspec}
\setmainfont{texgyrepagella}[Extension=.otf,UprightFont=*-regular,BoldFont=*-bold,ItalicFont=*-italic,BoldItalicFont=*-bolditalic]
\setsansfont{texgyreheros}[Extension=.otf,UprightFont=*-regular,BoldFont=*-bold,ItalicFont=*-italic,BoldItalicFont=*-bolditalic]
\setmonofont{texgyrecursor}[Extension=.otf,UprightFont=*-regular,BoldFont=*-bold,ItalicFont=*-italic,BoldItalicFont=*-bolditalic,Scale=MatchLowercase]
\usepackage{amsmath,amssymb}
\usepackage{unicode-math}
\setmathfont{texgyrepagella-math.otf}
\usepackage{xcolor,microtype,enumitem,needspace,fancyhdr}
\usepackage{bookmark}
\definecolor{ink}{HTML}{17344B}
\definecolor{accent}{HTML}{197D80}
\definecolor{muted}{HTML}{596773}
\hypersetup{colorlinks=true,linkcolor=accent,urlcolor=accent,pdftitle={RA-13 - Absolute-error
threshold for extremal Gaussian trace
bounds},pdfauthor={Open Problems in Numerical Linear Algebra}}
\urlstyle{same}
\setlength{\parindent}{0pt}
\setlength{\parskip}{5pt plus 1pt minus 1pt}
\linespread{1.04}
\setlength{\emergencystretch}{3em}
\widowpenalty=10000
\clubpenalty=10000
\displaywidowpenalty=10000
\allowdisplaybreaks[1]
\setlist{nosep,leftmargin=1.5em}
\providecommand{\tightlist}{\setlength{\itemsep}{0pt}\setlength{\parskip}{0pt}}
\providecommand{\pandocbounded}[1]{#1}
\pagestyle{fancy}
\fancyhf{}
\fancyhead[L]{\sffamily\footnotesize\color{muted}OPEN PROBLEMS IN NUMERICAL LINEAR ALGEBRA}
\fancyhead[R]{\sffamily\bfseries\footnotesize\color{accent}RA-13}
\fancyfoot[L]{\parbox[b]{0.9\textwidth}{\sffamily\fontsize{8}{10}\selectfont\color{muted}Editorial ratings; a bounded search cannot certify that a problem remains unsolved.\\Literature check: 2026-09-11}}
\fancyfoot[R]{\sffamily\footnotesize\color{muted}\thepage}
\renewcommand{\headrulewidth}{0.3pt}
\renewcommand{\headrule}{\hbox to\headwidth{\color{accent}\leaders\hrule height \headrulewidth\hfill}}
\makeatletter
\renewcommand{\section}{\Needspace{5\baselineskip}\@startsection{section}{1}{0pt}{12pt plus 2pt minus 2pt}{4pt}{\sffamily\bfseries\large\color{ink}}}
\renewcommand{\subsection}{\Needspace{4\baselineskip}\@startsection{subsection}{2}{0pt}{9pt plus 1pt minus 1pt}{3pt}{\sffamily\bfseries\normalsize\color{ink}}}
\makeatother
\setcounter{secnumdepth}{0}
\begin{document}
{\sffamily\small\color{muted}Randomized and low-rank approximation\par}
\vspace{3pt}
{\raggedright\hyphenpenalty=10000\exhyphenpenalty=10000\sffamily\bfseries\fontsize{21}{24}\selectfont\color{ink}Absolute-error
threshold for extremal Gaussian trace bounds\par}
\vspace{7pt}
{\sffamily\small\textbf{Difficulty:} challenging\quad\textbf{Importance:} interesting
to the community\par}
{\sffamily\small\bfseries\color{accent}Status: Solved\par}
\vspace{4pt}
{\color{accent}\hrule height 0.8pt}
\vspace{6pt}
\textbf{Rating rationale:} Challenging because signed eigenvalue
contributions require uniform extremal tail comparisons; community
impact is rigorous trace estimates for indefinite matrices.

\subsection{Resolution - 2026-09-11}\label{resolution---2026-09-11}

\textbf{Author:} George Stepaniants, Department of Computing and
Mathematical Sciences, California Institute of Technology, Pasadena,
California, USA.

\textbf{Solved affirmatively.} The
\href{https://github.com/ajt60gaibb/OpenProblemsInNLA/blob/main/randomized-and-low-rank-approximation/RA-13/solution.md}{complete
proof, Sections 1-8} establishes both probability comparisons for every
nonzero real symmetric matrix, including indefinite and zero-trace
matrices, every positive sample count, and the stated non-strict
threshold. Equations (1)-(2) give a stronger normalized one-sided
comparison; Section 8 converts it to the exact retained target.
\href{https://github.com/ajt60gaibb/OpenProblemsInNLA/blob/main/randomized-and-low-rank-approximation/RA-13/solution.pdf}{Proof
PDF} ·
\href{https://github.com/ajt60gaibb/OpenProblemsInNLA/blob/main/randomized-and-low-rank-approximation/RA-13/solution.tex}{Standalone
XeLaTeX source}.

A separate
\href{https://github.com/ajt60gaibb/OpenProblemsInNLA/blob/main/references/stepaniants-ra13-2026-09-11/verification/RA-13-independent-review.md}{independent
Codex-agent review} passed the complete proof and checked Kwaśnicki's
published bell-shape theorem and Hallman's coefficient derivative
directly. The new inflection identity and grouped transfers are supplied
in the manuscript. Substantial ChatGPT/Codex assistance is disclosed;
this is independent agent verification, not external human peer review
or formal certification.
\href{https://github.com/ajt60gaibb/OpenProblemsInNLA/blob/main/references/stepaniants-ra13-2026-09-11/README.md}{Submission
record and public-status audit}.

The original target and ID are retained below. The difficulty and
importance labels record the historical open question. The auxiliary
counterexamples below retain their original author and limited scope.

\subsection{Related auxiliary counterexamples ---
2026-09-11}\label{related-auxiliary-counterexamples-2026-09-11}

\textbf{Author:} Matthew J. Colbrook, Department of Applied Mathematics
and Theoretical Physics, University of Cambridge. \textbf{Independent
Codex-agent review: PASS for the limited result below.}

The centered augmented density has a genuine inflection point beyond the
proposed auxiliary upper bound. Together with the mode example, this
refutes the upper assertions of Hallman Conjectures 1 and 2.

\textbf{Historical scope:} This auxiliary result neither proved nor
refuted the complete absolute Gaussian trace-tail probability chain.
Failure of an auxiliary sufficient condition does not refute the final
tail comparisons. The full target is now resolved by the separate proof
above; the auxiliary counterexamples and their attribution remain
unchanged.

\textbf{Primary reference:}
\href{https://github.com/ajt60gaibb/OpenProblemsInNLA/blob/main/references/colbrook-transfer-2026-09-11/manuscripts/04_gamma_auxiliary_counterexamples.pdf}{complete
authored PDF},
\href{https://github.com/ajt60gaibb/OpenProblemsInNLA/blob/main/references/colbrook-transfer-2026-09-11/manuscripts/04_gamma_auxiliary_counterexamples.tex}{standalone
TeX}, \textbf{Proposition 3.1 (with Proposition 2.1 for related
evidence)}.
\href{https://github.com/ajt60gaibb/OpenProblemsInNLA/blob/main/references/colbrook-transfer-2026-09-11/verification/reviews/gamma-auxiliary-review.md}{Independent
proof review} ·
\href{https://github.com/ajt60gaibb/OpenProblemsInNLA/blob/main/references/colbrook-transfer-2026-09-11/README.md}{Authorship
and submission record}. Verification is independent agent review, not
external human peer review or formal certification.

\newpage

\subsection{Problem statement}\label{problem-statement}

For a real symmetric \(d\times d\) matrix \(D\), define

\[
T_m(D)=\frac1m\sum_{j=1}^{m}z_j^TDz_j,
\qquad z_1,\ldots,z_m\overset{\mathrm{iid}}{\sim}N(0,I_d).
\]

Let \(A\ne0\) be an arbitrary real symmetric \(n\times n\) matrix,
possibly indefinite. Put

\[
\lambda=\|A\|_2,\qquad
\phi=\|A\|_F,\qquad
\rho=\frac{\phi^2}{\lambda^2},\qquad
B_{\lambda,\phi}
=\lambda\operatorname{diag}\left(
I_{\lfloor\rho\rfloor},\sqrt{\rho-\lfloor\rho\rfloor}
\right).
\]

The norms are the spectral and Frobenius norms. Let \(X\) have Gamma
shape \(m\rho/2\) and rate \(m/(2\lambda)\), so that
\(\mathbb E X=\phi^2/\lambda\). Here a Gamma variable with shape
\(\alpha\) and rate \(\beta\) has density
\(\beta^\alpha x^{\alpha-1}e^{-\beta x}/\Gamma(\alpha)\) for \(x>0\).

\textbf{Conjecture.} For every integer \(n\ge1\), every such \(A\),
every integer \(m\ge1\), and every\\
\[
\varepsilon\ge
\frac{2\lambda}{m}
+\sqrt{\frac{2\phi^2}{m}
+\left(\frac{2\lambda}{m}\right)^2},
\]

the following comparisons hold:

\[
\begin{aligned}
\Pr\!\left(|T_m(A)-\operatorname{tr}(A)|\ge\varepsilon\right)
&\le
2\Pr\!\left(T_m(B_{\lambda,\phi})
-\operatorname{tr}(B_{\lambda,\phi})\ge\varepsilon\right)\\
&\le 2\Pr\!\left(X-\frac{\phi^2}{\lambda}\ge\varepsilon\right).
\end{aligned}
\]

Zero trailing diagonal entries in \(B_{\lambda,\phi}\) are harmless.
Each probability uses the appropriate estimator dimension. This is an
absolute-error question, even when \(\operatorname{tr}(A)=0\).

\subsection{Why it matters in numerical linear
algebra}\label{why-it-matters-in-numerical-linear-algebra}

The bound would quantify trace-estimation error even when positive and
negative eigenvalues cancel.

\subsection{References}\label{references}

\begin{enumerate}
\def\labelenumi{\arabic{enumi}.}
\tightlist
\item
  Eric Hallman,
  \href{https://arxiv.org/html/2411.15454v1\#S5}{\emph{Extremal bounds
  for Gaussian trace estimation}}, arXiv:2411.15454v1 (2024), §5,
  Theorem 7 and Conjecture 4.
\end{enumerate}

\subsection{Status check}\label{status-check}

Theorem 7 leaves its threshold unspecified. Checked 2026-09-08: the
source still lists only v1. Author, title and 2025/2026 searches,
including the \href{https://arxiv.org/abs/2512.02316}{later XTrace
paper}, found no resolution. This is a bounded check.

The displayed conjecture is retained despite a source prose
inconsistency: its threshold tends to zero as \(m\to\infty\), whereas
the following sentence says it tends to \(\phi\).

\subsection{Audit --- 2026-09-10}\label{audit-2026-09-10}

Rechecked \href{https://arxiv.org/html/2411.15454v1}{Hallman, Theorem 7
and Conjecture 4}; the record remains v1. Author and
extremal-Gaussian-trace searches found no resolution. The displayed
formula matches the numbered conjecture; its inconsistent following
prose remains disclosed, and no sharper threshold is claimed proved.
\end{document}
